\documentclass[11pt,a4paper]{article}
\usepackage[utf8]{inputenc}
\usepackage[T1]{fontenc}
\usepackage[spanish,es-noquoting,es-nodecimaldot]{babel}
\usepackage{amsmath,amssymb}
\usepackage{booktabs}
\usepackage{xcolor}
\usepackage[margin=2.3cm]{geometry}
\usepackage[colorlinks=true,linkcolor=blue!60!black]{hyperref}
\newcommand{\Ap}{A'}
\newcommand{\eps}{\epsilon}
\newcommand{\dd}{\mathrm{d}}
\title{Ampliación del informe: el umbral, el rango y la ecuación (50) del apunte}
\author{E.~L.~De Paoli\\[2pt]\small escrito con asistencia de agentes (Claude Opus 5)}
\date{21 de septiembre de 2026}
\begin{document}
\maketitle

\section{Qué es el umbral}
El umbral es la energía \emph{mínima del fotón incidente}, $E_\gamma$, para que $\gamma e^-\to \Ap e^-$ sea posible. No es la
energía máxima del fotón oscuro; esa es $x_+$, otra cantidad. En el centro de masa el mínimo es producir el $\Ap$ y el
electrón final en reposo relativo, $\sqrt s=m_e+M$. Con el electrón blanco en reposo, $s=m_e^2+2m_eE_\gamma$, así que
\begin{equation}
E_{\rm th}=M+\frac{M^2}{2m_e}.
\end{equation}
El término $M^2/2m_e$ es retroceso: el sistema final tiene que llevarse el impulso $\vec k$ del fotón, y eso cuesta energía
cinética que crece con $M$. Para $M=1$~MeV vale $0.98$~MeV: el umbral es $1.98$~MeV. Para $M=0.1$~MeV vale $0.01$~MeV y
no se nota.

\section{Dónde está el error en el apunte, y dónde no}
El apunte escribe ``umbral: $\omega\to m_{\Ap}$'' sólo en el Cuadro~1 de la sección 8.2. Ese cuadro no se usa en ninguna
ecuación que conduzca a la (50). Tu observación es correcta: \textbf{el umbral ingenuo del Cuadro~1 no entra en el
cálculo de (50) ni en el del número de fotones oscuros.}

El mismo error entra por otro lado: por el rango de integración, ecs.~(24)--(25),
\begin{equation}
\omega'\in\Big[\frac{m_e\omega+M^2/2}{2\omega+m_e},\;\omega\Big],
\label{eq:rango}
\end{equation}
que sí se usa cuando (50) se integra sobre $\omega'$ o se convoluciona con el espectro del reactor. El extremo superior
$\omega'_{\max}=\omega$ es el que está mal: para $M\neq0$ el $\Ap$ nunca se lleva toda la energía del fotón. El rango
exacto es
\begin{equation}
x_\pm=\frac{(E_\gamma+m_e)(s+M^2-m_e^2)\pm E_\gamma\sqrt{[s-(m_e+M)^2][s-(m_e-M)^2]}}{2s},
\end{equation}
y el umbral está adentro de esta fórmula: la raíz se anula exactamente en $E_\gamma=E_{\rm th}$, donde $x_+=x_-$. El
rango \eqref{eq:rango} del apunte, en cambio, es no vacío para todo $\omega\ge M/2$, y si además se pide $\omega'\ge M$
queda $\omega\ge M$: el umbral ingenuo del Cuadro~1 es lo que la ec.~\eqref{eq:rango} implica. Son el mismo error, escrito
dos veces; la copia del Cuadro~1 es inocua, la de la ec.~\eqref{eq:rango} no.

\section{Cuánto pesa cada error cuando se calcula el número de fotones oscuros}
Convolucioné la ec.~(50) del apunte, integrada sobre el rango \eqref{eq:rango} del apunte, con el espectro de Park a 1~GW y
$\eps=1$ (ecuación de fuente de Park con $\sigma_{\rm tot}=\sigma_{\rm KN}$), y lo mismo con la sección eficaz exacta sobre
$[x_-,x_+]$. Script: \texttt{work/check\_eq50\_fold.py}.

\begin{center}\small
\begin{tabular}{@{}lcccc@{}}
\toprule
$M$ [MeV] & fotones con $M<\omega<E_{\rm th}$ & \multicolumn{2}{c}{$\dd N/\dd E_{\Ap}$ apunte/exacto} & $N_{\Ap}$(3--8 MeV) apunte/exacto\\
 & & 3~MeV & 5~MeV & \\
\midrule
0.1 & 1.1\% & 1.06 & 1.05 & 1.05\\
0.5 & 23.6\% & 1.48 & 1.43 & 1.46\\
1.0 & 65.9\% & 2.05 & 1.90 & 1.99\\
\bottomrule
\end{tabular}
\end{center}

\noindent Lectura:
\begin{itemize}
\item En la ventana 3--8~MeV \textbf{el error de umbral/rango no pesa}. Un $\Ap$ de $E_{\Ap}\ge3$~MeV sólo puede venir de
un fotón con $\omega\ge3$~MeV, y todos esos están sobre el umbral para $M\le1$~MeV. El extremo superior $\omega'_{\max}=\omega$
en vez de $x_+$ cambia el límite inferior de $\omega$ de $3$ a $3.05$~MeV: nada.
\item Lo que pesa en la ventana es \textbf{la fórmula (50) misma}: factor $1.05$, $1.46$, $1.99$ para $M=0.1$, $0.5$,
$1$~MeV. Eso viene del $|\mathcal M|^2$ de la ec.~(40) (que descarta los términos en $m_e^2$) y del jacobiano sin masa.
Es el error que hay que corregir para usar (50) en TEXONO.
\item El error de rango pesa \textbf{a baja energía del $\Ap$}. Para $M=1$~MeV, el 66\% de los fotones por encima de $1$~MeV
está por debajo del umbral verdadero y el apunte los deja producir $\Ap$; además, para $E_{\Ap}$ entre $0.31$ y $1.34$~MeV la
(50) integra sobre una región cinemáticamente prohibida que contiene el polo $\omega'=M^2/2m_e=0.978$~MeV. En
$E_{\Ap}=1$~MeV el exacto da cero y el apunte da $1.5\times10^{21}$~MeV$^{-1}$s$^{-1}$. Para un detector con umbral bajo
(el PCGe de TEXONO, 0.3--12~keV de retroceso, que ve $\Ap$ de cualquier energía) esto sí importa.
\end{itemize}

\section{Corrección mínima para usar (50)}
Dos reemplazos y (50) queda bien: el corchete $[\ldots]$ por $F(a,b,c,d)/2$ con
$F=-2(a/b+b/a)+4(2c+d)[(1/a+1/b)+c(1/a+1/b)^2-d/(ab)]$, $a=2m_e\omega$, $b=M^2-2m_e\omega'$, $c=m_e^2$, $d=M^2$, y
el prefactor $\frac{\pi\eps^2\alpha^2}{m_e\omega^2}(1+\frac{M^2}{2m_e\omega})$ por $\frac{\pi\eps^2\alpha^2}{m_e\omega^2}$;
es decir, $\dd\sigma/\dd\omega'=\eps^2e^4F/(32\pi m_e\omega^2)$. Y el rango \eqref{eq:rango} por $[x_-,x_+]$. Con eso el
umbral sale solo y el polo queda afuera para todo $M<2m_e$.

\medskip\noindent\small\sloppy Recibo: 21/09/2026, \texttt{work/.venv/bin/python work/check\_eq50\_fold.py}; salida en
\texttt{work/check\_eq50\_fold.log} y \texttt{work/check\_eq50\_fold\_out.json}.
\end{document}
